\section{Metryki}
Jedną z metryk wykorzystanych w pracy było \gls{ASR}. Przy pomocy tej metryki mierzy się procent ataków na model, które z powodzeniem ominęły zabezpieczenia~\cite{correia2026systematicliteraturereviewllm}.
Metrykę tę definiuje się jako stosunek liczby udanych prób ataku do całkowitej liczby przeprowadzonych prób. Celem wprowadzenia mechanizmów obronnych była minimalizacja tej metryki. Może być zapisana wzorem~\eqref{eq:asr}.
\begin{equation}
    \label{eq:asr}
    \text{ASR}=\frac{\text{liczba udanych prób ataku}}{\text{liczba prób ataku}}\times100\%
\end{equation}

W celu oceny wpływu zaimplementowanych metod mitygacji podatności na wydajność systemu mierzono również opóźnienie.
Wnioskowanie AI dzieli się na fazę przetwarzania promptu oraz fazę sekwencyjnego generowania tokenów~\cite{agrawal2025evaluatingperformancellminference}.
Z tego względu latencję analizowano przy użyciu dwóch miar. Pierwszą z nich był czas do wygenerowania pierwszego tokenu (ang.~\gls{TTFT}), na który wpływ powinny mieć filtry wejściowe czy modele zewnętrzne. Zdefiniowana jest jako opóźnienie między otrzymaniem żądania a pojawieniem się pierwszego wyjściowego tokenu oraz obejmuje ono czas na planowanie i przetwarzanie zapytania~\cite{agrawal2025evaluatingperformancellminference}. Jest istotna dla aplikacji, które wymagają wysokiej responsywności~\cite{agrawal2025evaluatingperformancellminference}.
Drugą z kolei metryką jest czas pomiędzy kolejnymi tokenami (ang.~\gls{TBT}), który jest przydatny przy analizie wpływu wprowadzonych mechanizmów inżynierii aktywacji, a jej wielkość wpływa na odbieraną przez użytkownika prędkość modelu~\cite{agrawal2025evaluatingperformancellminference}.
By uniknąć maskowania przestojów w generowaniu odpowiedzi i w pełni ocenić wydajność, zgodnie z zaleceniami metodologicznymi~\cite{agrawal2025evaluatingperformancellminference}, nie opierano się na uśrednionym czasie, lecz wyniki interpretowano na podstawie rozkładów prawdopodobieństwa.